\documentclass[11pt]{article}

\usepackage{amsmath,amssymb,amsthm}

\newtheorem{lemma}{Lemma}

\title{Type Repeat Window Lemma}
\date{}

\begin{document}

\maketitle

\section*{Statement}

Let \(G\) be a graph of maximum degree \(\Delta\).

Let \(C_1,C_2\) be cycles whose vertex signatures agree under
FO\(^k\) radius-\(R\) neighborhood types.

Let

\[
D = C_1 \oplus C_2 .
\]

\begin{lemma}[Type Repeat Window]

There exists a constant \(u=u(k,\Delta,R)\) such that along any
agreement segment of the cycles, two vertices with identical
FO\(^k\) neighborhood type must occur within distance \(u\).

\end{lemma}

\section*{Idea}

The number of FO\(^k\) neighborhood types is finite:

\[
T = T(k,\Delta,R).
\]

Along any sufficiently long segment, the pigeonhole principle forces
two vertices to share the same FO\(^k\) type within a bounded window.

This repeated type forces local structural repetition, bounding the
length of agreement windows.

\end{document}
